% PAGES: 127-128

\noindent Eigenvalues and corresponding eigenvectors of a matrix.
\medskip

\noindent The characteristic polynomial of a matrix. Eigenvalues and the roots of the
characteristic polynomial. Multiplicity of eigenvalues.
\medskip

\noindent Linear independent eigenvectors. The number of eigenvalues and eigenvectors of a
matrix.
\medskip

\noindent Diagonal forms. Diagonalizable matrices. Matrices with a full set of linearly
independent eigenvectors.
\medskip

\noindent Diagonally dominant matrices. Gershgorin's Theorem. Nonsingularity of strictly
diagonally dominant matrices.
\medskip

\noindent Eigenvalues and eigenvectors of tridiagonal symmetric matrices.
\medskip

\section{Properties and identifying eigenvalues and eigenvectors}

\begin{definition}
If $A$ is a square matrix, then $\lambda$ is an eigenvalue of $A$ if and only if there
exists a column vector $v \neq 0$ such that
\begin{equation}
Av \;=\; \lambda v.
\end{equation}
The vector $v$ is called an eigenvector corresponding to $\lambda$.
\end{definition}

Note that, even if all the entries of the matrix $A$ are real numbers, an eigenvalue of
$A$ may be a complex number, and the corresponding eigenvector may have complex entries.

If $v$ is an eigenvector corresponding to the eigenvalue $\lambda$ of $A$, then, for any
constant $c$, the vector $cv$ is also an eigenvector of $A$ corresponding to $\lambda$,
since
\[
A(cv) \;=\; c(Av) \;=\; c(\lambda v) \;=\; \lambda(cv).
\]
However, $cv$ is not considered to be an eigenvector different than $v$. Only linearly
independent eigenvectors count as different eigenvectors. As seen in the example below,
it is possible for one eigenvalue to have more than one linearly independent eigenvector.

\begin{examplex}
This example is \textit{for illustration purposes only}; in practice, the eigenvalues
and the corresponding eigenvectors of a matrix are computed using numerical methods; see
section 4.4 for more details.

(i) Let
\[
A \;=\; \begin{pmatrix} -8 & 0 & 3\\ 3 & -2 & -1.5\\ -18 & 0 & 7\end{pmatrix}.
\]
We will show that\footnote{A simpler way to compute the eigenvalues and the corresponding
eigenvectors of the matrix $A$, using the characteristic polynomial of $A$, will be given
subsequently.} the matrix $A$ has eigenvalues $1$ and $-2$, with

\[
\text{eigenvectors} \quad \begin{pmatrix}0\\1\\0\end{pmatrix} \ \ \text{and} \ \
\begin{pmatrix}1\\0\\2\end{pmatrix} \quad \text{corresponding to} \quad -2
\]
\[
\text{eigenvector} \quad \begin{pmatrix}2\\-1\\6\end{pmatrix} \quad
\text{corresponding to} \quad 1.
\]

Denote by $\lambda$ an eigenvalue of $A$ and by $v = (v_i)_{i=1:3}$ a corresponding
eigenvector of $\lambda$. Then, $Av = \lambda v$, which can be written as the following
linear system:
\begin{align}
-8v_1 + 3v_3 &= \lambda v_1;\\
3v_1 - 2v_2 - 1.5v_3 &= \lambda v_2;\\
-18v_1 + 7v_3 &= \lambda v_3.
\end{align}

By eliminating $v_3$ from (4.2) and (4.4), we obtain that
\[
(\lambda^2 + \lambda - 2)v_1 \;=\; (\lambda+2)(\lambda-1)v_1 \;=\; 0.
\]
Thus, either $v_1 = 0$, or $\lambda = -2$, or $\lambda = 1$.

\begin{itemize}
\item If $v_1 = 0$, it follows from (4.2) that $v_3 = 0$. From (4.3), we find that
$-2v_2 = \lambda v_2$. Recall that $v \neq 0$, since $v$ is an eigenvector. Thus,
$v_2 \neq 0$, and therefore $\lambda = -2$. We conclude that $\lambda = -2$ is an
eigenvalue of $A$ with corresponding eigenvector
\[
\begin{pmatrix}0\\ v_2\\ 0\end{pmatrix} \;=\; v_2\begin{pmatrix}0\\1\\0\end{pmatrix}.
\]
Since multiplying of an eigenvector by a constant does not change the eigenvector, we
choose $v_2 = 1$ and conclude that $\begin{pmatrix}0\\1\\0\end{pmatrix}$ is an
eigenvector corresponding to the eigenvalue $\lambda = -2$.
\item If $\lambda = -2$, then (4.2) and (4.4) are both equivalent to $v_3 = 2v_1$, and
(4.3) holds true with no further constraint on $v_1$, $v_2$, or $v_3$. Thus, any
eigenvector corresponding to $\lambda = -2$ is of the form
\[
\begin{pmatrix}v_1\\ v_2\\ 2v_1\end{pmatrix} \;=\; v_1\begin{pmatrix}1\\0\\2\end{pmatrix}
+ v_2\begin{pmatrix}0\\1\\0\end{pmatrix}.
\]
% CONTINUES in sec04_c1_2.tex (page 129 / folio 113): this \item and the
% surrounding \begin{itemize} (the bulleted case analysis for the
% eigenvalue lambda = -2 of matrix A) are left OPEN here, and so is the
% enclosing \begin{examplex} opened above ("This example is for
% illustration purposes only..."). sec04_c1_2.tex must NOT reopen either
% environment -- it continues this item's prose directly ("We conclude
% that the eigenvalue lambda = -2 has two linearly independent
% eigenvectors..."), closes this \itemize once the lambda = 1 case for
% matrix A is done, then gives part (ii) (matrix B) and closes the
% \examplex with \end{examplex} at the printed square (page 129 / folio
% 113), immediately before Definition 4.2.
